\documentclass[11pt]{article}
\usepackage[margin=1in]{geometry}
\usepackage{booktabs}
\usepackage{amsmath}
\usepackage{amssymb}
\usepackage{enumitem}
\usepackage{tabularx}
\usepackage[hidelinks]{hyperref}
\newcommand{\pra}{\mathrm{PRA}}
\newcommand{\refhandle}{\mathrm{REF}}
\newcommand{\kv}{\mathrm{KV}}
\newcommand{\q}{\mathbf{Q}}
\newcommand{\kmat}{\mathbf{K}}
\newcommand{\vmat}{\mathbf{V}}

\setlength{\emergencystretch}{3em}

\title{Standalone Progressive Retrieval Attention: A PyTorch Architecture for URI-Addressed Memory Expansion}
\author{Jorge Simao\\EInnovator R\&D -- Center for the Sciences of Computation, Cognition, and Complexity}
\date{\today}

\begin{document}
\maketitle

\begin{abstract}
This paper describes a standalone PyTorch research prototype for Progressive Retrieval Attention (PRA). The prototype is intentionally small: a decoder-only transformer is augmented with reference handles, an in-memory URI resolver, a per-layer key/value memory cache, and a PRA attention branch that performs cross-attention over selected reference memories. The system is designed for controlled experiments on synthetic, hierarchical, code, documentation, Wikipedia, book, and GitHub-style datasets where each question carries explicit reference metadata and expected anchors. We formalize the model, document the implemented cache-construction path, define the training and evaluation pipeline, and specify ablations required to separate ordinary language-model capacity from reference-memory use. The present manuscript is implementation-backed but not yet a report of completed experiments. All numerical result tables are mock placeholders and must be replaced by measured results before publication as an empirical paper.
\end{abstract}

\section{Introduction}

Progressive Retrieval Attention (PRA) is motivated by a common mismatch in long-context language modeling. Many tasks require access to large external context, but the relevant evidence is sparse, hierarchical, and often discovered only after partial reasoning. A conventional decoder-only transformer receives a flat token sequence. Retrieval-augmented generation improves this interface by prepending retrieved passages, but the retrieved text still becomes ordinary prompt context \cite{lewis2020rag}. PRA instead represents external context through lightweight reference handles that can be resolved into URI-addressed memory and attended to through a dedicated memory path.

This paper documents the first standalone PyTorch prototype in the PRA program. The prototype is not a production long-context model. It is a controlled experimental instrument. Its goals are to make the model definition inspectable, expose reference use through explicit metadata, and support ablations that would be difficult to interpret in a large pretrained model.

The implementation consists of four main components. First, a tiny GPT-style language model provides causal decoder blocks. Second, a reference system parses tokens such as \texttt{<REF\_1>} and maps them to URI-addressed references. Third, a resolver and cache encode referenced documents separately, storing per-layer key/value tensors. Fourth, a PRA attention layer retrieves cache entries by summary-vector similarity and adds a scaled cross-attention memory branch to ordinary causal self-attention.

\section{Model Definition}

\subsection{Decoder Backbone}

Let $x_{1:T}$ be a token sequence. The model embeds tokens and absolute positions:

\begin{equation}
H^{(0)} = E_{\mathrm{tok}}(x_{1:T}) + E_{\mathrm{pos}}(1:T).
\end{equation}

For each layer $\ell \in \{0,\ldots,L-1\}$, the implemented block is a pre-normalized residual decoder block:

\begin{align}
\tilde{H}^{(\ell)} &= \mathrm{LN}_1(H^{(\ell)}),\\
A^{(\ell)} &= \mathrm{PRAttn}_{\ell}(\tilde{H}^{(\ell)}),\\
U^{(\ell)} &= H^{(\ell)} + A^{(\ell)},\\
H^{(\ell+1)} &= U^{(\ell)} + \mathrm{FFN}_{\ell}(\mathrm{LN}_2(U^{(\ell)})).
\end{align}

The feed-forward network is a two-layer MLP with GELU activation and hidden width $4d$. The output head maps the final normalized hidden state to vocabulary logits. This matches the implementation in \texttt{TinyPRAModel} and \texttt{PRATransformerBlock}.

\subsection{Causal Self-Attention}

For hidden states $X \in \mathbb{R}^{B \times T \times d}$, the local branch computes multi-head causal self-attention:

\begin{align}
Q &= XW_Q,\quad K = XW_K,\quad V = XW_V,\\
O_{\mathrm{local}} &= W_O \,
\mathrm{merge}\left[
\mathrm{softmax}\left(
\frac{QK^\top + M_{\mathrm{causal}}}{\sqrt{d_h}}
\right)V
\right].
\end{align}

The implementation stores a lower-triangular causal mask with maximum sequence length configured by \texttt{PRAConfig.max\_seq\_len}.

\section{Reference Handles and Resolver}

\subsection{Reference Tokens}

The prototype uses explicit textual handles of the form \texttt{<REF\_n>}. A regular expression parser extracts occurrences and optionally attaches entries from a \texttt{ReferenceTable}. Each reference table entry stores an id, token, URI, summary, and metadata. Legacy \texttt{!!ref:uri!!} syntax is retained only for compatibility.

\subsection{URI and Anchor Semantics}

References point to simple URIs such as \path{mem://doc1}. Anchored locations use URI fragments, for example \path{mem://doc1#section.a}. Utility functions split the URI base from the anchor, compute parent anchors, and construct child anchors. The current in-memory resolver stores a dictionary from URI to text plus an optional dictionary from URI to summary. If a URI is requested with \texttt{summary\_only=True}, the resolver returns the summary in place of the full text. Child references are exposed when stored keys extend the requested URI with hierarchical suffixes.

\subsection{Resolved References}

Given a handle $r$ with URI $u$, the resolver returns

\begin{equation}
\mathrm{resolve}(u) = (\mathrm{text}_u, \mathrm{summary}_u, \mathrm{children}_u, \mathrm{metadata}_u).
\end{equation}

This simple resolver is sufficient for synthetic and local experiments. Future systems can replace it with file, database, vector-index, search, or web resolvers while preserving the same high-level interface.

\section{Layer-Specific Memory Cache}

\subsection{Cache Entries}

The PRA cache stores entries of the form

\begin{equation}
C_u = (u,\; \mathrm{text}_u,\; \mathrm{summary}_u,\; s_u,\; \{(K_{u,\ell},V_{u,\ell})\}_{\ell}).
\end{equation}

Here $s_u \in \mathbb{R}^{d}$ is a summary vector used for retrieval, and $(K_{u,\ell},V_{u,\ell})$ are detached key/value tensors projected for layer $\ell$. The implementation exposes an abstract \texttt{PRAMemoryCache} and a dictionary-backed \texttt{PRASimpleMemoryCache}.

\subsection{Reference Encoding Pass}

Before a batch forward pass, the training code builds a cache from batch metadata. For each referenced document:

\begin{enumerate}[leftmargin=*]
\item Resolve the URI to text and summary.
\item Tokenize the reference text and truncate to \texttt{max\_seq\_len}.
\item Tokenize the summary and compute a prototype summary vector by averaging summary-token embeddings.
\item Run the reference text through the same model path with PRA memory disabled.
\item At each layer, take the pre-attention normalized hidden state and project it through that layer's $W_K$ and $W_V$.
\item Store the resulting detached layer-specific tensors in the cache.
\end{enumerate}

This design is deliberately conservative. It avoids recursive PRA during first-level cache construction and ensures that external memory for layer $\ell$ uses the same projection matrices as the layer that will consume it.

\section{PRA Attention}

\subsection{Summary-Based Retrieval}

At generation or training time, if the cache is non-empty, the current implementation uses the last hidden state before attention as a retrieval query. For a query $q \in \mathbb{R}^{d}$ and cache summary vectors $s_u$, retrieval scores are cosine similarities:

\begin{equation}
\mathrm{score}(u \mid q) =
\frac{q^\top s_u}{\lVert q\rVert_2 \lVert s_u\rVert_2}.
\end{equation}

The cache returns the top-$k$ entries. Entries with similarity below \texttt{trigger\_threshold} are ignored. The current default hyperparameters are \texttt{top\_k\_refs=2} and \texttt{trigger\_threshold=0.2}.

\subsection{Memory Cross-Attention Branch}

For selected URI set $\mathcal{U}$ and layer $\ell$, the implementation concatenates cached keys and values across the selected entries:

\begin{equation}
K_{\mathcal{U},\ell} = \mathrm{concat}_{u \in \mathcal{U}} K_{u,\ell},
\quad
V_{\mathcal{U},\ell} = \mathrm{concat}_{u \in \mathcal{U}} V_{u,\ell}.
\end{equation}

It then computes a memory branch using the same query tensor $Q$ from the local branch:

\begin{equation}
O_{\mathrm{mem}} =
W_{O,\mathrm{mem}}
\mathrm{merge}\left[
\mathrm{softmax}\left(
\frac{QK_{\mathcal{U},\ell}^{\top}}{\sqrt{d_h}}
\right)V_{\mathcal{U},\ell}
\right].
\end{equation}

The returned attention output is

\begin{equation}
O = O_{\mathrm{local}} + \alpha O_{\mathrm{mem}},
\end{equation}

where $\alpha$ is \texttt{memory\_alpha}, defaulting to $0.5$. If the cache is empty, PRA is disabled, no selected entry passes the threshold, or the selected entry lacks layer-$\ell$ memory, the layer returns only the local self-attention output.

\subsection{Cross-Attention Rather than Native KV Injection}

The standalone prototype uses a separate cross-attention memory branch. Direct concatenation into native self-attention KV is listed as an ablation and future integration path, but it is intentionally not the default mechanism. Cross-attention cleanly separates local causal attention from external memory attention and avoids immediate complications from rotary position embeddings, grouped-query attention, and production cache layout. This makes the prototype suitable for controlled research before adapting large pretrained models.

\section{Dataset and DataLoader Pipeline}

\subsection{Dataset Schema}

Datasets are represented by three JSONL files per stage:

\begin{itemize}[leftmargin=*]
\item \texttt{documents.jsonl}: URI-addressed documents, summaries, titles, and anchors.
\item \texttt{references.jsonl}: reference ids, URIs, summaries, anchors, and metadata.
\item \texttt{questions.jsonl}: prompts, answers, expected reference ids, and expected anchors.
\end{itemize}

The loaded sample type is \texttt{QuestionSample}, containing a question, answer, list of \texttt{ReferenceSample} objects, target reference ids, and metadata. The current registry includes synthetic memory QA, hierarchical synthetic references, code repositories, Wikipedia, books, technical documentation, and GitHub-style repositories.

\subsection{Collation}

The \texttt{PRACollator} tokenizes each question and answer, forms next-token labels, pads variable-length rows, builds a per-sample reference table, and preserves non-tensor metadata. This metadata is essential: the training step uses it to resolve and encode references before the main model forward pass.

\subsection{Data Module}

The \texttt{PRADataModule} instantiates a dataset stage, builds a character-level tokenizer from questions, answers, summaries, and reference texts, constructs a collator, and creates train/validation/test splits. This pipeline keeps the prototype independent from pretrained tokenizers while still exercising the full reference-memory path.

\section{Training Objective and System}

The training objective is standard autoregressive cross-entropy:

\begin{equation}
\mathcal{L} =
-\sum_{t=1}^{T}
\log p_{\theta}(x_{t+1}\mid x_{\leq t}, C),
\end{equation}

where $C$ is the batch-specific PRA cache constructed from metadata. The training loop uses AdamW, optional learning-rate warmup, gradient accumulation, gradient clipping, optional mixed precision, checkpointing, early stopping, and TensorBoard/W\&B/ClearML-compatible logging hooks.

The implemented PRA batch step is:

\begin{enumerate}[leftmargin=*]
\item Move tensor fields to the target device.
\item Build the PRA cache from batch metadata.
\item Attach the cache to the model and all attention blocks.
\item Run the language-model forward pass with PRA memory enabled.
\item Compute cross-entropy with padding ignored.
\end{enumerate}

\section{Evaluation Metrics}

The evaluation code reports both language-model and reference-behavior metrics:

\begin{itemize}[leftmargin=*]
\item \textbf{Loss and perplexity}: autoregressive language-model quality.
\item \textbf{Answer accuracy}: exact containment of the expected answer in generated output or decoded continuation.
\item \textbf{Reference retrieval accuracy}: whether expected reference URIs appear in the selected cache.
\item \textbf{Expected anchor hit}: whether expected anchors appear in cached text or selected URI strings.
\item \textbf{Expanded reference count}: average number of expanded references per example.
\item \textbf{Average retrieved tokens}: tokenized size of resolved reference text.
\item \textbf{Full-context token count}: token cost of a full-context baseline.
\item \textbf{Cache hit ratio}: fraction of examples whose expected references are represented in cache.
\item \textbf{Latency and throughput}: wall-clock evaluation time, examples per second, and tokens per second.
\end{itemize}

These metrics are designed to test the PRA hypothesis directly. A useful PRA model should not merely achieve high answer accuracy; it should do so while selecting the right references and using fewer retrieved tokens than a full-context baseline.

\section{Experimental Protocol}

The intended baseline variants are:

\begin{enumerate}[leftmargin=*]
\item \textbf{No references}: question only, PRA cache cleared.
\item \textbf{Full context}: all available reference text prepended to the question.
\item \textbf{Simple RAG}: reference summaries prepended to the question.
\item \textbf{PRA}: question with reference handles and batch-specific PRA cache.
\end{enumerate}

The required ablations are:

\begin{itemize}[leftmargin=*]
\item No PRA memory branch.
\item Summary-only reference memory.
\item One-level retrieval.
\item Recursive retrieval through child anchors.
\item Cross-attention memory branch versus native KV injection.
\item Varying PRA-enabled layers.
\item Varying \texttt{top\_k\_refs}, \texttt{trigger\_threshold}, and \texttt{memory\_alpha}.
\end{itemize}

The current implementation exposes configuration fields for several of these ablations. Some, such as native KV injection and full recursive learned expansion, remain future work and should be marked as such in experimental reports.

\section{Mock Results Template}

Table~\ref{tab:mock-results} is a placeholder showing the intended reporting format. The values are illustrative mock numbers only.

\begin{table}[h]
\centering
\caption{Mock result table. Values are illustrative placeholders and are not experimental measurements.}
\label{tab:mock-results}
\begin{tabular}{lrrrrr}
\toprule
Variant & Accuracy & Ref hit & Tokens in & Retrieved & Latency \\
\midrule
No references & 0.00 & 0.00 & 96 & 0 & 1.00x \\
Full context & 0.00 & 1.00 & 512 & 512 & 1.80x \\
Simple RAG & 0.00 & 0.60 & 160 & 64 & 1.15x \\
PRA & 0.00 & 0.00 & 96 & 0 & 1.00x \\
\bottomrule
\end{tabular}
\end{table}

Before publication as an empirical paper, this table must be regenerated from real runs and accompanied by configuration files, seeds, checkpoint identifiers, dataset versions, and confidence intervals or repeated-run variance where feasible.

\section{Relationship to Prior Work}

The prototype builds on the transformer attention architecture \cite{vaswani2017attention}. It is related to sparse and long-context attention models such as Longformer, BigBird, Reformer, and Routing Transformer \cite{beltagy2020longformer,zaheer2020bigbird,kitaev2020reformer,roy2021routing}, but the central question differs: PRA studies whether external context can remain addressable until needed rather than being placed in a single long sequence.

It also relates to memory-augmented and retrieval-conditioned models. Compressive Transformers maintain compressed memories over earlier activations \cite{rae2020compressive}. Memorizing Transformers retrieve nearest-neighbor memories \cite{wu2022memorizing}. RETRO conditions generation on retrieved chunks from a large database \cite{borgeaud2022retro}. RAG and REALM connect generation to non-parametric retrieval \cite{lewis2020rag,guu2020realm}. PRA differs in its emphasis on explicit handles, URI-addressed references, recursive anchors, and per-layer cache construction.

Finally, the implementation anticipates runtime concerns raised by serving systems such as vLLM's PagedAttention \cite{kwon2023vllm}. A mature PRA system would require cache admission, sharing, eviction, batching, and memory safety policies.

\section{Limitations}

The current prototype has important limitations. The tokenizer is character-level and intended for controlled experiments rather than competitive language modeling. Summary vectors are computed by averaging summary-token embeddings, which is simple but likely weaker than learned summary encoders. Retrieval currently uses the last hidden state rather than explicit attention to reference-token positions. Recursive anchor expansion is represented in the resolver interface but is not yet a fully learned iterative policy. Evaluation contains containment-style answer checks that are useful for synthetic QA but insufficient for open-ended generation. Finally, no real experimental superiority is claimed in this manuscript.

\section{Conclusion}

The standalone PyTorch PRA prototype provides a concrete testbed for demand-driven, URI-addressed memory expansion. Its main contribution is not a final architecture but an inspectable experimental substrate: explicit reference handles, resolver-backed memory, per-layer K/V cache construction, a cross-attention memory branch, and metrics that measure reference behavior alongside answer quality. This substrate enables the next step in the PRA program: controlled experiments that determine when progressive retrieval attention improves the tradeoff among accuracy, context cost, latency, and interpretability.

\bibliographystyle{plain}
\bibliography{../shared/references}
\end{document}
